% GENERATED FILE. Edit canonical lessons, then run npm run generate:books.
% canonical-source-hash: fnv1a64:d0a19f7f154584b6
% canonical-lessons: GE-C451-reduzieren, GE-C451-sonderangebot, GE-C451-schnaeppchen, GE-C451-feilschen, GE-C451-flohmarkt

\chapter{Gute Angebote}
\label{ch:ge-gute-angebote}
\hlchaptermodality{451}

\begin{chapteropening}
\textbf{By the end of this chapter:} \emph{I can find a bargain: reducing prices, the special offer, the bargain, haggling and the flea market.}

Complete the last lesson of chapter 451: I can find a bargain: reducing prices, the special offer, the bargain, haggling and the flea market.
\end{chapteropening}

\section[reduzieren]{reduzieren — to reduce}

\label{lesson:GE-C451-reduzieren}



\begin{quote}
Before the new one: say the German for to redeem, then the German for to pay out.
\end{quote}

\subsection*{You'll want to know: reduzieren}
\textbf{reduzieren} — \textquotedblleft{}to reduce\textquotedblright{}.

\begin{quote}
Im Januar reduziert das Kaufhaus alle Preise.
\end{quote}

\begin{quote}
\emph{In January the department store reduces all its prices.}
\end{quote}

Past participle without \emph{ge-}: \emph{reduziert} — reduced, on price tags.

\begin{cousinweb}
\emph{reduzieren} is a learned borrowing from Latin \emph{redūcere} 'to lead back, to bring back' (Lewis \& Short); EtymDB lists the German word.
\end{cousinweb}

\begin{grammarlens}[title={What you've built}]
One more word for telling, explaining and describing at B1.
\end{grammarlens}

\subsection*{Your turn}
\begin{itemize}
\raggedright
  \item \emph{Say it:} \emph{reduzieren}
  \item \emph{Say it:} \emph{reduzieren}, once more
  \item \emph{Say it:} a sentence of your own with \emph{reduzieren}
  \item \emph{Recall:} say the German for to redeem, then the German for to pay out, then say \emph{reduzieren} again
\end{itemize}

\begin{tcolorbox}[breakable,colback=teal!4,colframe=teal!35!black,title={Before you move on}]
What does reduzieren mean? (\textquotedblleft{}To reduce\textquotedblright{}.) Say it once more.
\end{tcolorbox}
\section[das Sonderangebot]{das Sonderangebot — the special offer}

\label{lesson:GE-C451-sonderangebot}



\begin{quote}
Before the new one: say the German for to pay out, then the German for to reduce.
\end{quote}

\subsection*{You'll want to know: das Sonderangebot}
\textbf{das Sonderangebot} — \textquotedblleft{}the special offer\textquotedblright{}.

\begin{quote}
Die Erdbeeren sind diese Woche im Sonderangebot.
\end{quote}

\begin{quote}
\emph{Strawberries are on special offer this week.}
\end{quote}

Plural: \emph{die Sonderangebote}.

\begin{cousinweb}
\emph{Sonderangebot} is built from \emph{sonder-} 'special, apart', inherited from Proto-Germanic \emph{sundraz} (a cousin of English \emph{asunder}), and \emph{Angebot} 'offer', from \emph{anbieten} 'to offer'.
\end{cousinweb}

\begin{grammarlens}[title={What you've built}]
One more word for telling, explaining and describing at B1.
\end{grammarlens}

\subsection*{Your turn}
\begin{itemize}
\raggedright
  \item \emph{Say it:} \emph{das Sonderangebot}
  \item \emph{Say it:} \emph{das Sonderangebot}, once more
  \item \emph{Say it:} a sentence of your own with \emph{das Sonderangebot}
  \item \emph{Recall:} say the German for to pay out, then the German for to reduce, then say \emph{das Sonderangebot} again
\end{itemize}

\begin{tcolorbox}[breakable,colback=teal!4,colframe=teal!35!black,title={Before you move on}]
What does das Sonderangebot mean? (\textquotedblleft{}The special offer\textquotedblright{}.) Say it once more.
\end{tcolorbox}
\section[das Schnäppchen]{das Schnäppchen — the bargain}

\label{lesson:GE-C451-schnaeppchen}



\begin{quote}
Before the new one: say the German for to reduce, then the German for the special offer.
\end{quote}

\subsection*{You'll want to know: das Schnäppchen}
\textbf{das Schnäppchen} — \textquotedblleft{}the bargain\textquotedblright{}.

\begin{quote}
Die Jacke war ein Schnäppchen: nur zwanzig Euro.
\end{quote}

\begin{quote}
\emph{The jacket was a bargain: only twenty euros.}
\end{quote}

Same form in the plural. \emph{Ein Schnäppchen machen} — to get a bargain.

\begin{cousinweb}
\emph{Schnäppchen} is built on \emph{schnappen} 'to snap up, to grab' (EtymDB lists English \emph{snap} as a cousin) with \emph{-chen} 'little': something quickly snapped up.
\end{cousinweb}

\begin{grammarlens}[title={What you've built}]
One more word for telling, explaining and describing at B1.
\end{grammarlens}

\subsection*{Your turn}
\begin{itemize}
\raggedright
  \item \emph{Say it:} \emph{das Schnäppchen}
  \item \emph{Say it:} \emph{das Schnäppchen}, once more
  \item \emph{Say it:} a sentence of your own with \emph{das Schnäppchen}
  \item \emph{Recall:} say the German for to reduce, then the German for the special offer, then say \emph{das Schnäppchen} again
\end{itemize}

\begin{tcolorbox}[breakable,colback=teal!4,colframe=teal!35!black,title={Before you move on}]
What does das Schnäppchen mean? (\textquotedblleft{}The bargain\textquotedblright{}.) Say it once more.
\end{tcolorbox}
\section[feilschen]{feilschen — to haggle}

\label{lesson:GE-C451-feilschen}



\begin{quote}
Before the new one: say the German for the special offer, then the German for the bargain.
\end{quote}

\subsection*{You'll want to know: feilschen}
\textbf{feilschen} — \textquotedblleft{}to haggle\textquotedblright{}.

\begin{quote}
Auf dem Markt feilscht er um jeden Euro.
\end{quote}

\begin{quote}
\emph{At the market he haggles over every euro.}
\end{quote}

With \emph{um}: \emph{um den Preis feilschen}.

\begin{cousinweb}
\emph{feilschen} is usually linked to the old adjective \emph{feil} 'for sale', but the details are not settled, and EtymDB does not list the word.
\end{cousinweb}

\begin{grammarlens}[title={What you've built}]
One more word for telling, explaining and describing at B1.
\end{grammarlens}

\subsection*{Your turn}
\begin{itemize}
\raggedright
  \item \emph{Say it:} \emph{feilschen}
  \item \emph{Say it:} \emph{feilschen}, once more
  \item \emph{Say it:} a sentence of your own with \emph{feilschen}
  \item \emph{Recall:} say the German for the special offer, then the German for the bargain, then say \emph{feilschen} again
\end{itemize}

\begin{tcolorbox}[breakable,colback=teal!4,colframe=teal!35!black,title={Before you move on}]
What does feilschen mean? (\textquotedblleft{}To haggle\textquotedblright{}.) Say it once more.
\end{tcolorbox}
\section[der Flohmarkt]{der Flohmarkt — the flea market}

\label{lesson:GE-C451-flohmarkt}



\begin{quote}
Before the new one: say the German for the bargain, then the German for to haggle.
\end{quote}

\subsection*{You'll want to know: der Flohmarkt}
\textbf{der Flohmarkt} — \textquotedblleft{}the flea market\textquotedblright{}.

\begin{quote}
Am Sonntag gibt es einen Flohmarkt am Fluss.
\end{quote}

\begin{quote}
\emph{On Sunday there is a flea market by the river.}
\end{quote}

Plural: \emph{die Flohmärkte}.

\begin{cousinweb}
\emph{Flohmarkt} is built from \emph{Floh} 'flea', inherited from Proto-Germanic \emph{flauhaz} and a cousin of English \emph{flea}, and \emph{Markt} 'market', from Latin \emph{mercātus} (EtymDB).
\end{cousinweb}

\begin{grammarlens}[title={What you've built}]
One more word for telling, explaining and describing at B1.
\end{grammarlens}

\subsection*{Your turn}
\begin{itemize}
\raggedright
  \item \emph{Say it:} \emph{der Flohmarkt}
  \item \emph{Say it:} \emph{der Flohmarkt}, once more
  \item \emph{Say it:} a sentence of your own with \emph{der Flohmarkt}
  \item \emph{Recall:} say the German for the bargain, then the German for to haggle, then say \emph{der Flohmarkt} again
\end{itemize}

\begin{tcolorbox}[breakable,colback=teal!4,colframe=teal!35!black,title={Before you move on}]
What does der Flohmarkt mean? (\textquotedblleft{}The flea market\textquotedblright{}.) Say it once more.
\end{tcolorbox}
